\documentclass[12pt]{article}
\usepackage{latexblog}

% ============================================================
% Blog Metadata
% ============================================================
\blogtitle{Netty(1)：介绍}
\blogdate{2020-09-02}
\blogtags{框架, 闲话编程}
\bloglang{zh}
\blogsummary{}

\begin{document}
\makeblogtitle

Netty是一个高性能的异步事件驱动的网络应用框架，本质上是对NIO进行了高层的抽象，使得可以轻松的创建服务器和客户端，极大简化了诸如TCP和UDP套接字的操作。

\section{入门}\label{ux5165ux95e8}

\subsection{核心概念}\label{ux6838ux5fc3ux6982ux5ff5}

\begin{figure}
\centering
\pandocbounded{\includegraphics[keepaspectratio,alt={Netty architecture}]{https://netty.io/3.8/guide/images/architecture.png}}
\caption{Netty architecture}
\end{figure}

Netty中的一些核心概念：

\begin{itemize}
\tightlist
\item
  \texttt{Channel}对应到Socket
\item
  \texttt{EventLoop}用来控制流、多线程处理以及并发
\item
  \texttt{ChannelFuture}用来实现异步通知
\end{itemize}

\subsection{ECHO server}\label{echo-server}

最简单的例子是构建一个echo server，发过来什么同样返回什么。

首先需要实现一个Handler，定义如何处理消息：

\begin{Shaded}
\begin{Highlighting}[]
\KeywordTok{public} \KeywordTok{class}\NormalTok{ EchoServerHandler }\KeywordTok{extends}\NormalTok{ ChannelInboundHandlerAdapter }\OperatorTok{\{}
    \AttributeTok{@Override}
    \KeywordTok{public} \DataTypeTok{void} \FunctionTok{channelRead}\OperatorTok{(}\NormalTok{ChannelHandlerContext ctx}\OperatorTok{,} \BuiltInTok{Object}\NormalTok{ msg}\OperatorTok{)} \KeywordTok{throws} \BuiltInTok{Exception} \OperatorTok{\{}
        \CommentTok{// handler需要去释放msg对象(引用计数）}
        \CommentTok{// 这里不用去手动release msg，因为writeAndFlush里面已经处理了}
\NormalTok{        ctx}\OperatorTok{.}\FunctionTok{writeAndFlush}\OperatorTok{(}\NormalTok{msg}\OperatorTok{);}
    \OperatorTok{\}}

    \AttributeTok{@Override}
    \KeywordTok{public} \DataTypeTok{void} \FunctionTok{exceptionCaught}\OperatorTok{(}\NormalTok{ChannelHandlerContext ctx}\OperatorTok{,} \BuiltInTok{Throwable}\NormalTok{ cause}\OperatorTok{)} \KeywordTok{throws} \BuiltInTok{Exception} \OperatorTok{\{}
\NormalTok{        cause}\OperatorTok{.}\FunctionTok{printStackTrace}\OperatorTok{();}
\NormalTok{        ctx}\OperatorTok{.}\FunctionTok{close}\OperatorTok{();}
    \OperatorTok{\}}
\OperatorTok{\}}
\end{Highlighting}
\end{Shaded}

然后利用\texttt{ServiceBootStrap}来启动

\begin{Shaded}
\begin{Highlighting}[]
\KeywordTok{public} \KeywordTok{class}\NormalTok{ EchoServer }\OperatorTok{\{}
    \KeywordTok{private} \DataTypeTok{final} \DataTypeTok{int}\NormalTok{ port}\OperatorTok{;}

    \KeywordTok{public} \FunctionTok{EchoServer}\OperatorTok{(}\DataTypeTok{int}\NormalTok{ port}\OperatorTok{)} \OperatorTok{\{}
        \KeywordTok{this}\OperatorTok{.}\FunctionTok{port} \OperatorTok{=}\NormalTok{ port}\OperatorTok{;}
    \OperatorTok{\}}

    \KeywordTok{public} \DataTypeTok{void} \FunctionTok{run}\OperatorTok{()} \KeywordTok{throws} \BuiltInTok{InterruptedException} \OperatorTok{\{}
\NormalTok{        EventLoopGroup bossGroup }\OperatorTok{=} \KeywordTok{new} \FunctionTok{NioEventLoopGroup}\OperatorTok{();}
\NormalTok{        EventLoopGroup workerGroup }\OperatorTok{=} \KeywordTok{new} \FunctionTok{NioEventLoopGroup}\OperatorTok{();}
        \ControlFlowTok{try} \OperatorTok{\{}
\NormalTok{            ServerBootstrap bootstrap }\OperatorTok{=} \KeywordTok{new} \FunctionTok{ServerBootstrap}\OperatorTok{();}
\NormalTok{            bootstrap}\OperatorTok{.}\FunctionTok{group}\OperatorTok{(}\NormalTok{bossGroup}\OperatorTok{,}\NormalTok{ workerGroup}\OperatorTok{)}
                    \OperatorTok{.}\FunctionTok{channel}\OperatorTok{(}\NormalTok{NioServerSocketChannel}\OperatorTok{.}\FunctionTok{class}\OperatorTok{)}
                    \OperatorTok{.}\FunctionTok{childHandler}\OperatorTok{(}\KeywordTok{new}\NormalTok{ ChannelInitializer}\OperatorTok{\textless{}}\BuiltInTok{SocketChannel}\OperatorTok{\textgreater{}()} \OperatorTok{\{}
                        \AttributeTok{@Override}
                        \KeywordTok{protected} \DataTypeTok{void} \FunctionTok{initChannel}\OperatorTok{(}\BuiltInTok{SocketChannel}\NormalTok{ ch}\OperatorTok{)} \KeywordTok{throws} \BuiltInTok{Exception} \OperatorTok{\{}
\NormalTok{                            ch}\OperatorTok{.}\FunctionTok{pipeline}\OperatorTok{().}\FunctionTok{addLast}\OperatorTok{(}\KeywordTok{new} \FunctionTok{EchoServerHandler}\OperatorTok{());}
                        \OperatorTok{\}}
                    \OperatorTok{\})}
                    \OperatorTok{.}\FunctionTok{option}\OperatorTok{(}\NormalTok{ChannelOption}\OperatorTok{.}\FunctionTok{SO\_BACKLOG}\OperatorTok{,} \DecValTok{128}\OperatorTok{)}
                    \OperatorTok{.}\FunctionTok{childOption}\OperatorTok{(}\NormalTok{ChannelOption}\OperatorTok{.}\FunctionTok{SO\_KEEPALIVE}\OperatorTok{,} \KeywordTok{true}\OperatorTok{);}

\NormalTok{            ChannelFuture f }\OperatorTok{=}\NormalTok{ bootstrap}\OperatorTok{.}\FunctionTok{bind}\OperatorTok{(}\NormalTok{port}\OperatorTok{).}\FunctionTok{sync}\OperatorTok{();}
\NormalTok{            f}\OperatorTok{.}\FunctionTok{channel}\OperatorTok{().}\FunctionTok{closeFuture}\OperatorTok{().}\FunctionTok{sync}\OperatorTok{();}
        \OperatorTok{\}} \ControlFlowTok{finally} \OperatorTok{\{}
\NormalTok{            workerGroup}\OperatorTok{.}\FunctionTok{shutdownGracefully}\OperatorTok{();}
\NormalTok{            bossGroup}\OperatorTok{.}\FunctionTok{shutdownGracefully}\OperatorTok{();}
        \OperatorTok{\}}
    \OperatorTok{\}}

    \KeywordTok{public} \DataTypeTok{static} \DataTypeTok{void} \FunctionTok{main}\OperatorTok{(}\BuiltInTok{String}\OperatorTok{[]}\NormalTok{ args}\OperatorTok{)} \KeywordTok{throws} \BuiltInTok{InterruptedException} \OperatorTok{\{}
        \KeywordTok{new} \FunctionTok{EchoServer}\OperatorTok{(}\DecValTok{9999}\OperatorTok{).}\FunctionTok{run}\OperatorTok{();}
    \OperatorTok{\}}
\OperatorTok{\}}
\end{Highlighting}
\end{Shaded}

\section{基本原理}\label{ux57faux672cux539fux7406}

\subsection{EventLoop与Server Channel、Channel}\label{eventloopux4e0eserver-channelchannel}

其运行原理如图：

\begin{figure}
\centering
\pandocbounded{\includegraphics[keepaspectratio,alt={EventLoopGroup and Channels}]{https://dpzbhybb2pdcj.cloudfront.net/maurer/Figures/03fig04_alt.jpg}}
\caption{EventLoopGroup and Channels}
\end{figure}

\begin{itemize}
\tightlist
\item
  左边只有一个ServerChannel，代表服务器上监听某个端口的套接字，所以实际上也只需要一个EventLoop就可以了
\item
  右边代表建立的客户端连接，每一个连接都对应到一个EventLoop，当有很多个连接的时候，这些连接是会共享其中的EventLoop的。
\end{itemize}

\subsection{ChannelPipeline}\label{channelpipeline}

每次建立连接的时候，都会调用\texttt{ChannelInitializer}，这个类负责安装一些自定义的\texttt{ChannelHandler}到\texttt{ChannelPipeline}中。 实际上的程序可能对应到多个入站和出站的ChannelHandler，它们的执行顺序是由它们被添加的顺序所决定的，类似这样：

\begin{figure}
\centering
\pandocbounded{\includegraphics[keepaspectratio,alt={Channel Pipeline}]{https://dpzbhybb2pdcj.cloudfront.net/maurer/Figures/03fig03_alt.jpg}}
\caption{Channel Pipeline}
\end{figure}

\subsection{ChannelHandler}\label{channelhandler}

ChanelHandler又有很多的类型，比如：

\begin{itemize}
\tightlist
\item
  编码器和解码器, 例如\texttt{ByteToMessageDecoder}、\texttt{ProtobufEncoder}等
\item
  \texttt{SimpleChannelInboundHandler\textless{}T\textgreater{}}，用来处理简单的逻辑比如收到消息后完成业务逻辑，只需要实现其中的\texttt{void\ channelRead0(ChannelHandlerContext\ ctx,\ I\ msg)}方法即可
\end{itemize}

\subsection{Bootstrap}\label{bootstrap}

前面使用\texttt{ServerBootstrap}类来启动了服务器上的socket监听，如果是客户端程序可以使用\texttt{Bootstrap}类来完成。一个比较明显的区别就是，客户端程序只需要一个\texttt{EventLoopGroup}而服务端通常会需要两个。

Ref:

\begin{itemize}
\tightlist
\item
  \href{https://netty.io/wiki/user-guide-for-4.x.html}{Netty User guide for 4.x}
\item
  \href{https://livebook.manning.com/book/netty-in-action}{Netty in Action}
\end{itemize}


\end{document}
